
%\needspace{4\baselineskip}
\color{uniblue}{\subsection{Максимальная токовая защита трансформатора (МТЗ)}}
\color{black}

\begin{enumerate}[label=\arabic{section}.\arabic{subsection}.\arabic{enumi}, labelsep=4pt, leftmargin=0pt, itemindent=57pt, itemsep=0pt, parsep=5pt] % если что можно крутить itemsep
%%%%%%%%%%%%%%%%%%%%%%%%%%%%%%%%%%%%%%%%%%%%%%%%%%%%%%%%%%%  ОБЩИЕ СВЕДЕНИЯ %%%%%%%%%%%%%%%%%%%%%%%%%%%%%%%%%%%%%%%%%%%%%%%%

\item Общие сведения

\begin{enumerate}[label=\arabic{section}.\arabic{subsection}.\arabic{enumi}.\arabic*, labelsep=4pt, leftmargin=0em, itemindent=65pt, parsep=0pt]

\item
Максимальная токовая защита применяется для защиты обмоток трансформатора от сверхтоков внешних КЗ при повреждениях на присоединениях прилегающей сети, а также для резервирования основных защит трансформатора и подключается к трансформаторам тока стороны ВН трансформатора. МТЗ имеет возможность работы с контролем напряжения стороны НН трансформатора.

\item
В устройстве максимальная токовая защита представлена функциональным блоком «МТЗ». В составе функционального блока предусмотрены следующие функции:
\begin{itemize}
\item три ступени максимальной токовой защиты (МТЗ 1 ст, МТЗ 2 ст, МТЗ 3 ст);
\item комбинированный пусковой орган напряжения (КПОН);
\item блокировка при броске намагничивающего тока (БНТ);
\item блокировка логической защиты шин (БЛЗШ).
\end{itemize}

\item
В целях предотвращения действия ступеней МТЗ при коротких замыканиях в сетях 6-35 кВ (в случаях, когда нейтраль трансформатора заземлена) предусматривается возможность  исключения токов нулевой последовательности посредством переключения измеряемых токов с фазных на линейные, для этого программный переключатель SGF1 «Сборка\_ток\_цепей» переводится в состояние «Треугольник».

\item
МТЗ может быть полностью выведена из работы оперативно путем активации сигнала «ОВ МТЗ».

\end{enumerate}

%%%%%%%%%%%%%%%%%%%%%%%%%%%%%%%%%%%%%%%%%%%%%%%%%%%%%%%%%  СТУПЕНЬ МТЗ %%%%%%%%%%%%%%%%%%%%%%%%%%%%%%%%%%%%%%%%%%%%%%%%

\item Ступени максимальной токовой защиты (МТЗ 1 ст, МТЗ 2 ст, МТЗ 3 ст)

\begin{enumerate}[label=\arabic{section}.\arabic{subsection}.\arabic{enumi}.\arabic*, labelsep=4pt, leftmargin=0em, itemindent=65pt, parsep=0pt]

\item
В функциональном блоке «МТЗ» ступени МТЗ представлены функциями «МТЗ 1 ст», «МТЗ 2 ст», «МТЗ 3 ст».

\item
Ввод каждой ступени МТЗ в работу на этапе параметрирования устройства осуществляется своим программным переключателем SGF1 «Ввод\_функции». Информация о введенной ступени отображается выходным сигналом для каждой ступени («МТЗ 1 ст: Ввод», «МТЗ 2 ст: Ввод», «МТЗ 3 ст: Ввод»).

\item
Каждая из ступеней МТЗ может быть выведена из работы оперативно путем активации сигнала оперативного вывода («ОВ МТЗ 1 ст», «ОВ МТЗ 2 ст», «ОВ МТЗ 3 ст»), что характеризуется активным состоянием сигнала «Оперативный вывод» («МТЗ 1 ст: Оперативный вывод», «МТЗ 2 ст: Оперативный вывод», «МТЗ 3 ст: Оперативный вывод») на выходе функции.

\item
При активном состоянии входного сигнала «МТЗ 1 ст на сигн», «МТЗ 2 ст на сигн», «МТЗ 3 ст на сигн» действие соответствующей ступени МТЗ переводится на сигнал.

\item
Ступень МТЗ контролирует повышение любого фазного тока ($I_A, I_B, I_C$) или линейного тока ($I_{AB}, I_{BC}, I_{CA}$) (в зависимости от положения программного переключателя SGF1 «Сборка\_ток\_цепей») выше заданной уставки, при превышении заданной уставки по току через выдержку времени таймера T1 (уставка «Тср») выдается сигнал на отключение трансформатора. 

\item\sloppy
Для ступени МТЗ возможны два режима работы с контролем напряжения секции НН: «Вольтметровая блокировка» и «Управляющее напряжение» (выбор режима работы осуществляется программным переключателем SGF2 «Тип\_КПОН»). Для выбора режима работы без контроля напряжения секции НН необходимо перевести программный переключатель SGF5 «Реж\_КПОН» в состояние «Не предусмотрено».

\item
В режиме «Вольтметровая блокировка» при превышении уставки по току для срабатывания через выдержку времени таймера T1 (уставка «Тср») необходимо подтверждение от КПОН (снижение напряжения на секции НН). При неисправностях в цепях напряжения ступень МТЗ может или выводиться из работы, или действовать на отключение, в зависимости от выбранного состояния программного переключателя SGF4 «Реж\_БНН\_КПОН»: «Деблокировка» или «Блокировка». 

\item
Для работы в режиме «Управляющее напряжение» необходимо задать две уставки по току для ступени МТЗ: «Iср» -- уставка по току срабатывания и «Iср\_груб» -- уставка по ток срабатывания грубого органа. В указанном режиме при отсутствии неисправности цепей напряжения при превышении уставки по току «Iср» для срабатывания необходимо подтверждение от КПОН (снижение напряжения на секции НН). При неисправностях в цепях напряжения и при превышении уставки по току грубого органа «Iср\_груб» через выдержку времени таймера T1 (уставка T1 «Тср») выдается сигнал срабатывания. Режим работы на уставке «Iср\_груб» определяется состоянием программного переключателя SGF4 «Реж\_БНН\_КПОН»: «Деблокировка» или «Блокировка».

\item
При необходимости возможна блокировка выбранных ступеней МТЗ при броске тока намагничивания силового трансформатора от БНТ, это достигается путем перевода программного переключателя выбора режима контроля от БНТ SGF3 «Реж\_БНТ» в положение «Предусмотрено».

\item
Для ступеней МТЗ предусмотрена возможность отслеживания положения СВ стороны НН. Для правильной работы в этом режиме необходимо выделить две ступени МТЗ и выставить программный переключатель SGF6 «Контр\_СВ» в положение «Блокировка ступени при включенном СВ» для ступени с ускоренной выдержкой времени и в положение «Блокировка ступени при отключенном СВ» для ступени с нормальной выдержкой времени. В этом случае, при изменении состояния СВ НН, автоматически вводится в работу одна из ступеней со своими уставками, учитывающими изменение режима работы сети.

\item
Алгоритм работы ступеней МТЗ выполнен в соответствии с рисунком \ref{mtzvn:img1}. В таблице \ref{mtzvn:tbl1} приведены параметры, необходимые для настройки МТЗ.


\vspace{3mm}
\begin{figure}[!h]
\centering
\includegraphics[width=1\textwidth,height=1\textheight,keepaspectratio]{\fbpath/05. МТЗ Т 2обм/_latex/img/mtzst.pdf}
\captionof{figure}{Функциональная схема логики ступеней МТЗ}\label{mtzvn:img1}
\end{figure}

\end{enumerate}
\FloatBarrier % Все float'ы ДО этой строки будут размещены здесь.
%%%%%%%%%%%%%%%%%%%%%%%%%%%%%%%%%%%%%%%%%%%%%%%%%%%%%%%%% ОРГАН КПОН %%%%%%%%%%%%%%%%%%%%%%%%%%%%%%%%%%%%%%%%%%%%%%%%

%\color{uniblue}{\subsubsection{Комбинированный пусковой орган напряжения (КПОН)}}
%\color{black}

\item Комбинированный пусковой орган напряжения (КПОН)

\begin{enumerate}[label=\arabic{section}.\arabic{subsection}.\arabic{enumi}.\arabic*, labelsep=4pt, leftmargin=0em, itemindent=65pt, parsep=0pt]

\item
В функциональном блоке «МТЗ» комбинированный пусковой орган напряжения представлен функцией «КПОН».

\item
Ввод функции осуществляется автоматически при активации режим пуска от внутреннего КПОН введенной в работу ступени МТЗ.

\item
Положением программного переключателя SGF1 «Реж\_пуска» осуществляется выбор алгоритма работы функции:
\begin{itemize}
\item по снижению минимального из трех линейных напряжений $U_{AB}, U_{BC}, U_{CA}$ ниже уставки $U_{\textsl{мин}}$ (состояние программного переключателя «По Uмин»);
\item комбинированный: по снижению минимального из трех линейных напряжений $U_{AB}, U_{BC}, U_{CA}$ ниже уставки $U_{\textsl{мин}}$ или по повышению напряжения обратной последовательности $U_2$ выше уставки $U_{2\textsl{макс}}$ (состояние программного переключателя «Комбинированный»);
\item при приеме сигнала от внешнего КПОН «КПОН внеш.» (состояние программного переключателя «Внешний»).
\end{itemize}

\item
Функция контролирует снижение любого из линейных напряжений ($U_{AB}, U_{BC}, U_{CA}$) ниже заданной уставки или, при введенном положении «Комбинированный» программного переключателя «SGF1», повышение напряжения обратной последовательности $U_2$ выше заданной уставки или прием сигнала «КПОН внеш.», в случае выполнения заданных условий выдается сигнал пуска КПОН.

\item
В качестве воздействующих величин функция использует действующие значения основной гармоники линейных напряжений и напряжение обратной последовательности.

\item
Алгоритм работы функции выполнен в соответствии с рисунком \ref{mtzvn:img2}.

\vspace{3mm}
\begin{figure}[H]
\centering
\includegraphics[width=0.7\textwidth,height=0.7\textheight,keepaspectratio]{\fbpath/05. МТЗ Т 2обм/_latex/img/kpon.pdf}
\captionof{figure}{Функциональная схема логики КПОН}\label{mtzvn:img2}
\end{figure}

\end{enumerate}
\FloatBarrier % Все float'ы ДО этой строки будут размещены здесь.

%%%%%%%%%%%%%%%%%%%%%%%%%%%%%%%%%%%%%%%%%%%%%%%%%%%%%%%%% ОРГАН БЛЗШ %%%%%%%%%%%%%%%%%%%%%%%%%%%%%%%%%%%%%%%%%%%%%%%%

%\color{uniblue}{\subsubsection{Блокировка логической защиты шин (БЛЗШ)}}
%\color{black}
\item Блокировка логической защиты шин (БЛЗШ)

\begin{enumerate}[label=\arabic{section}.\arabic{subsection}.\arabic{enumi}.\arabic*, labelsep=4pt, leftmargin=0em, itemindent=65pt, parsep=0pt]

\item
Функция блокировки логической защиты шин (БЛЗШ) предусматривается для блокировки ЛЗШ на вводном и секционном выключателях секции, питающей силовой трансформатор. В качестве воздействующих величин используются сигналы пуска ступеней МТЗ. Состояние программного переключателя SGF1 «БлокЛЗШ\_выбор\_ст» определяет ступень МТЗ, от которой будет производиться блокировка ЛЗШ.

\item
В функциональном блоке «МТЗ» блокировка логической защиты шин представлена функцией «БЛЗШ».

\item
Алгоритм работы функции выполнен в соответствии с рисунком \ref{mtzvn:img4}.

\begin{figure}[H]
\centering
\includegraphics[width=0.6\textwidth,height=0.6\textheight,keepaspectratio]{\fbpath/05. МТЗ Т 2обм/_latex/img/blzsh.pdf}
\captionof{figure}{Функциональная схема логики БЛЗШ}\label{mtzvn:img4}
\vspace{3mm}
\end{figure}

\end{enumerate}


%%%%%%%%%%%%%%%%%%%%%%%%%%%%%%%%%%%%%%%%%%%%%%%%%%%%%%%%% ОРГАН БНТ %%%%%%%%%%%%%%%%%%%%%%%%%%%%%%%%%%%%%%%%%%%%%%%%

%\color{uniblue}{\subsubsection{Блокировка при броске намагничивающего тока (БНТ)}}
%\color{black}
\item Блокировка при броске намагничивающего тока (БНТ)

\begin{enumerate}[label=\arabic{section}.\arabic{subsection}.\arabic{enumi}.\arabic*, labelsep=4pt, leftmargin=0em, itemindent=65pt, parsep=0pt]

\item
Функция блокировки при броске намагничивающего тока позволяет заблокировать защиту при броске тока намагничивания, возникающем при постановке силового трансформатора под напряжение. 

\item
В функциональном блоке «МТЗ» блокировка при броске намагничивающего тока представлена функцией «БНТ». Ввод функции осуществляется автоматически при активации режим контроля броска тока намагничивания введенной в работу ступени МТЗ.

\item
Условие пуска выполняется при одновременном наличии следующих условий:
\begin{itemize}
\item отношение действующих значений фазных токов второй гармоники ($I_{A\ 2\textsl{г}}, I_{B\ 2\textsl{г}}, I_{C\ 2\textsl{г}}$) к действующим значениям фазных токов основной гармоники ($I_A, I_B, I_C$) превышает уставку «Ih2/Ih1» (контролируется для каждой фазы отдельно);
\item отсутствует срабатывание порогового элемента, реагирующего на превышение максимального из действующих значений основной гармоники фазных токов ($I_A, I_B, I_C$) над уставкой «I\_ср»;
\item сформирован пуск одной из ступеней МТЗ (контроль выполняется пофазный).
\end{itemize}

\item
В функции предусмотрена возможность ввода перекрестной блокировки: при переводе программного переключателя SGF1 «Перекрест\_блок» в состояние «Предусмотрено» обнаружение броска тока намагничивания в одной из фаз блокирует срабатывание по всем фазам ступени МТЗ. Пуск БНТ учитывается отдельно для каждой из ступеней МТЗ при помощи соответствующего программного переключателя «Реж\_БНТ» в составе ступени МТЗ. При этом блокировка каждой ступени осуществляется пофазно.
\item
Алгоритм работы функции выполнен в соответствии с рисунком \ref{mtzvn:img3}.

\vspace{2mm}
\begin{figure}[H]
\centering
\includegraphics[width=1.0\textwidth,height=1.0\textheight,keepaspectratio]{\fbpath/05. МТЗ Т 2обм/_latex/img/bnt.pdf}
\captionof{figure}{Функциональная схема логики БНТ}\label{mtzvn:img3}
\end{figure}

\end{enumerate}
\FloatBarrier % Все float'ы ДО этой строки будут размещены здесь.


\end{enumerate}
\FloatBarrier % Все float'ы ДО этой строки будут размещены здесь.